\section{The two merge deficits on a common copy space}
\label{sec:pair_merge_deficit}

Let $Q,C,V,R$ be finite orthonormal bases, with $m\ge0$ copies of each.
Use the same enumeration of the copies in all four factors, and put
\begin{align}
 \mathcal H_{QC}&=\ell^2(Q)^{\otimes m}\otimes\ell^2(C)^{\otimes m},\notag\\
 \mathcal H_{VR}&=\ell^2(V)^{\otimes m}\otimes\ell^2(R)^{\otimes m}.\notag
\end{align}
The common space is $\mathcal H_{QC}\otimes\mathcal H_{VR}$.
The separate and simultaneous copy actions of
Definition~\ref{def:pair_copy_actions} define the actual central
logarithmic dimension observables in each pair. These are the two local
merge deficits used in \cite[Section~7]{OpenAI2026AreaLaw}.

\begin{definition}[The actual paired merge deficit]
\label{def:pair_merge_deficit}
\lean{TensorPower.pairMergeDeficit,TensorPower.isHermitian_pairMergeDeficit}
\leanok
\uses{def:pair_copy_actions}
For the $QC$ pair set $D_C=F_Q+F_C-F_{QC}$, and for the $VR$ pair set
$D_R=F_V+F_R-F_{VR}$. Each operator is Hermitian. On the common space
write $\widehat D_C=D_C\otimes I_{VR}$ and
$\widehat D_R=I_{QC}\otimes D_R$.
\end{definition}
\begin{proof}
\leanok
Each logarithmic dimension observable is a real linear combination of
orthogonal central projectors. Sums and differences of Hermitian
operators remain Hermitian, as do their extensions by an identity.
\end{proof}

\begin{theorem}[The two lifted exponential factors]
\label{thm:pair_merge_deficit_lifts}
\lean{TensorPower.commute_pairMergeDeficit_lifts,TensorPower.exp_pairMergeDeficit_lifts}
\leanok
\uses{def:pair_merge_deficit}
The lifted deficits commute. For all real $a,b$,
\begin{align}
 e^{a\widehat D_C+b\widehat D_R}
 &=e^{aD_C}\otimes e^{bD_R}.
 \label{eq:pair_merge_deficit_lifts}
\end{align}
\end{theorem}
\begin{proof}
\leanok
The two products $\widehat D_C\widehat D_R$ and
$\widehat D_R\widehat D_C$ both equal $D_C\otimes D_R$.
For commuting operators the exponential of their sum is the product of
their exponentials. Exponentiating either tensor-factor extension
adjoins the same identity, giving (\ref{eq:pair_merge_deficit_lifts}).
\end{proof}

\begin{theorem}[The individual traces on the common space]
\label{thm:pair_merge_deficit_traces}
\lean{TensorPower.trace_exp_pairMergeDeficit_left,TensorPower.trace_exp_pairMergeDeficit_right}
\leanok
\uses{def:pair_merge_deficit}
Let $\rho$ be any matrix on the common space, and let
$\rho_{QC}=\operatorname{Tr}_{VR}\rho$ and
$\rho_{VR}=\operatorname{Tr}_{QC}\rho$. For all real $a,b$,
\begin{align}
 \Tr(\rho e^{a\widehat D_C})&=\Tr(\rho_{QC}e^{aD_C}),
 \label{eq:pair_merge_deficit_trace_left}\\
 \Tr(\rho e^{b\widehat D_R})&=\Tr(\rho_{VR}e^{bD_R}).
 \label{eq:pair_merge_deficit_trace_right}
\end{align}
The matrix need not be normalized, invariant under copy permutations,
or independent between the two pairs. The identities are complex trace
equalities, and hence also hold for their real parts.
\end{theorem}
\begin{proof}
\leanok
The partial trace is adjoint to adjoining an identity under the trace
pairing. Apply this identity to each local exponential, using the
identity for its tensor-factor extension.
\end{proof}

% The component construction must still identify these two reduced
% matrices with the actual local component marginals.
A joint exponential trace retains the correlations in $\rho$; no product of marginal trace
pairings follows from the operator factorization alone.
